\section{Funcional HMT--MD para la anomalía magnética del muón}
\label{sec:anomalia-magnetica-muon-rev9}

Para un leptón de factor giromagnético \(g_\ell\), la anomalía magnética
adimensional se define por
\[
 a_\ell:=\frac{g_\ell-2}{2},
 \qquad
 g_\ell=2(1+a_\ell).
\]
En el caso muónico, la construcción HMT--MD utiliza como coordenadas de
entrada la constante de estructura fina interna \(\alpha_{\rm HMT}\), la
autoescala \(\varphi=(1+\sqrt5)/2\), la completación decimal de base mil y
el defecto aritmético \(54\) de APP. Estas cantidades preceden a la
evaluación del observable y pertenecen a dominios ya construidos en la
primera parte del tratado.

Sea
\[
 \mathcal C_\mu
 :=\{(\alpha,\varphi)\in\mathbb R_{>0}^2\}.
\]
Sobre este espacio se define el funcional adimensional
\begin{equation}
 \mathfrak a_\mu(\alpha,\varphi)
 :=\frac{\alpha}{2\pi}
 +\frac{\alpha(\varphi-1)}{1000}
 +\frac{\alpha^2}{1000(54+1)}.
 \label{eq:funcional-g2-muon-rev9}
\end{equation}
Los tres sumandos mantienen funciones diferenciadas. El primero es el canal
de Schwinger, normalizado por el cierre angular \(2\pi\). El segundo acopla la
autoescala \(\varphi-1=\varphi^{-1}\) con la completación decimal
\(1000=729+243+27+1\). El tercero es el cierre cuadrático de \(\alpha\)
sobre el sucesor \(54+1\) del defecto APP. La ecuación
\eqref{eq:funcional-g2-muon-rev9} es, por construcción, un escalar
adimensional y su evaluación precede a cualquier comparación experimental.

\begin{proposition}[Evaluación interna del funcional muónico]
Al especializar \eqref{eq:funcional-g2-muon-rev9} en
\[
 (\alpha,\varphi)
 =\left(\alpha_{\rm HMT},\frac{1+\sqrt5}{2}\right),
\]
se obtiene
\[
 \boxed{\mathfrak a_\mu^{\rm HMT}
 =0.0011659207130080142},
\]
y, por la definición del factor giromagnético,
\[
 \boxed{\mathfrak g_\mu^{\rm HMT}
 =2\bigl(1+\mathfrak a_\mu^{\rm HMT}\bigr)
 =2.002331841426016}.
\]
\end{proposition}

\begin{proof}
La primera igualdad resulta de la sustitución directa de
\(\alpha_{\rm HMT}\) y \(\varphi\) en
\eqref{eq:funcional-g2-muon-rev9}. La segunda aplica la identidad
\(g_\mu=2(1+a_\mu)\). Ambas operaciones se realizan en
\(\mathbb R_{>0}\) y conservan el carácter adimensional del observable.
\end{proof}

\section{Contraste metrológico posterior}
\label{sec:contraste-g2-muon-rev9}

La referencia adoptada para el contraste es el promedio de las campañas
Run--1 a Run--6 del experimento Muon \(g-2\) de Fermilab
\cite{MuonGMinus2_2025}:
\[
 a_\mu^{\rm ref}=0.001165920705,
 \qquad
 u(a_\mu^{\rm ref})=1.48\times10^{-10}.
\]
La diferencia entre la evaluación HMT--MD y esa referencia es
\[
 \mathfrak a_\mu^{\rm HMT}-a_\mu^{\rm ref}
 =8.0080142\times10^{-12},
\]
equivalente a
\[
 \frac{\mathfrak a_\mu^{\rm HMT}-a_\mu^{\rm ref}}
 {u(a_\mu^{\rm ref})}
 =0.054108204\ldots.
\]
La referencia y su incertidumbre intervienen exclusivamente en este cociente
de contraste. Una actualización experimental modifica la última evaluación y
conserva inalterado el funcional interno
\eqref{eq:funcional-g2-muon-rev9}.

\paragraph{Estatuto físico y dominio de validez del funcional leptónico.}
El teorema establecido en esta sección comprende la definición tipada del
funcional, su evaluación exacta una vez fijadas
\(\alpha_{\rm HMT}\) y \(\varphi\), la conversión a
\(\mathfrak g_\mu\) y el contraste metrológico posterior. La selección
estructural conjunta de los tres correctores de
\eqref{eq:funcional-g2-muon-rev9} a partir de un Hamiltoniano local requiere
un teorema adicional. Esta delimitación conserva la evaluación demostrada y
fija con precisión el contenido de la afirmación física.
